%% ============================================================
%%  paper2_credit_equivalence.tex  (v2 — March 2026)
%%  "From Perpetual Contracts to Islamic Credit:
%%   The Random Stopping Time Equivalence and
%%   a Riba-Free Foundation for Sukuk, Takaful, and Credit Protection"
%%  Shehzad Ahmed — February 2026 (revised March 2026)
%% ============================================================
\documentclass[10pt,twocolumn]{article}
\usepackage[utf8]{inputenc}
\usepackage[margin=0.75in]{geometry}
\usepackage{amsmath,amssymb,amsthm}
\usepackage{graphicx}
\usepackage{hyperref}
\usepackage[numbers]{natbib}
\usepackage{booktabs}
\usepackage{abstract}
\usepackage{xurl}
\usepackage{titlesec}
\usepackage{authblk}
\usepackage{xcolor}
\usepackage{array}
\usepackage[protrusion=true,expansion=false]{microtype}
\usepackage{enumitem}
\usepackage{caption}
\usepackage{multirow}
\usepackage{bm}

%% Section spacing
\titlespacing*{\section}{0pt}{13pt}{6pt}
\titlespacing*{\subsection}{0pt}{10pt}{4pt}
\titlespacing*{\subsubsection}{0pt}{8pt}{3pt}

%% Abstract style
\renewcommand{\abstractnamefont}{\normalfont\bfseries\large}
\renewcommand{\abstracttextfont}{\normalfont\small}

%% Theorem environments
\theoremstyle{plain}
\newtheorem{theorem}{Theorem}
\newtheorem{proposition}{Proposition}
\newtheorem{corollary}{Corollary}
\newtheorem{lemma}{Lemma}
\theoremstyle{definition}
\newtheorem{definition}{Definition}
\newtheorem{remark}{Remark}
\newtheorem{example}{Example}

%% Math operators
\DeclareMathOperator{\E}{\mathbb{E}}
\newcommand{\Qa}{Q^a}
\newcommand{\R}{\mathbb{R}}
\newcommand{\Prob}{\mathbb{P}}
\newcommand{\F}{\mathcal{F}}

\definecolor{ribared}{RGB}{180,30,30}
\definecolor{halalgreen}{RGB}{30,120,60}
\definecolor{noteblue}{RGB}{20,60,160}

\hypersetup{colorlinks=true,linkcolor=blue,citecolor=blue,urlcolor=blue}

%% ============================================================
\title{\LARGE\textbf{From Perpetual Contracts to Islamic Credit:\\
The Random Stopping Time Equivalence and\\
a Riba-Free Foundation for Sukuk,\\
Takaful, and Credit Protection}}

\author[1]{Shehzad Ahmed}

\affil[1]{Arcus Quant Fund; Baraka Protocol\\
\small\texttt{shehzad@arcusquantfund.com}}

\date{February 2026 (revised March 2026)}
%% ============================================================

\begin{document}

\twocolumn[\begin{@twocolumnfalse}
\maketitle\thispagestyle{empty}

\begin{abstract}
\noindent
We establish a formal mathematical equivalence between two previously unconnected
pricing frameworks --- the random time representation of perpetual futures
contracts \cite{Ackerer2024} and the reduced-form credit risk model of Duffie
\& Singleton \cite{DuffieSingleton1999} --- and demonstrate that this
equivalence has profound implications for Islamic finance.

The core observation is structural: the stopping time $\theta_t$ in the
Ackerer-Hugonnier-Jermann perpetual pricing formula is isomorphic to the
default time $\tau$ in the Duffie-Singleton defaultable bond formula,
under the identification of the convergence intensity $\kappa$ with the hazard
rate $\lambda$ and the interest parameter $\iota$ with the risk-free
discount rate $r$. When the interest parameter is set to zero ($\iota = 0$)
--- which Ackerer et al.\ prove achieves unique no-arbitrage
perpetual futures pricing --- the analogous credit instrument at $r = 0$
is also no-arbitrage, with the stopping time distribution alone carrying
all timing information previously encoded in the interest discount.

This equivalence resolves a 50-year impasse in Islamic finance: it provides
a mathematically rigorous, no-arbitrage pricing foundation for credit
instruments --- sukuk, takaful, and credit default swaps --- that requires
no positive interest rate. We prove the main theorem, derive closed-form
prices for an Islamic perpetual sukuk and an everlasting takaful contract
using Proposition~6 of Ackerer et al., and show that the convergence
intensity $\kappa$ functions as a riba-free analog to the risk-free rate
across all Islamic credit instruments. We term this the \textbf{$\kappa$-rate}.
We further establish a \textbf{Separation Theorem}: interest is a
mathematically separable, eliminable component of no-arbitrage pricing ---
not an intrinsic necessity. Empirical evidence from the European negative
interest rate environment (2014--2022) corroborates the theoretical finding
that credit markets function without a positive discount rate.

\medskip
\noindent\textbf{Keywords:} \textit{riba}, Islamic finance, credit risk,
sukuk, takaful, perpetual contracts, stopping time, no-arbitrage,
Duffie-Singleton, $\kappa$-rate, riba-free pricing, blockchain sukuk

\medskip
\noindent\textbf{JEL Codes:} G12, G13, G15, G22, Z12
\end{abstract}
\bigskip
\end{@twocolumnfalse}]

%% ============================================================
\section{Introduction}
%% ============================================================

The pricing of credit risk is the central unresolved problem in Islamic
finance. The \$3 trillion Islamic finance industry \cite{IFSB2024} operates
without any mathematical tool for pricing credit default risk that does not
require a positive interest rate. Every credit derivatives framework in use
today --- the structural models of Merton \cite{Merton1974}, Black \&
Cox \cite{BlackCox1976}, and Leland \cite{Leland1994}; the reduced-form
models of Jarrow \& Turnbull \cite{JarrowTurnbull1995}, Duffie \&
Singleton \cite{DuffieSingleton1999,DuffieSingleton2003}, and Lando
\cite{Lando1998}; the textbook synthesis of Bielecki \& Rutkowski
\cite{BieleckiRutkowski2002} and Jeanblanc, Yor \& Chesney
\cite{JeanblanYorChesney2009} --- is built on a positive risk-free
interest rate $r > 0$. Islamic scholars correctly identify this $r$ as
\textit{riba}: it is the compensation for the pure time value of money,
which Islamic jurisprudence prohibits as a predetermined excess in exchange
transactions \cite{ElGamal2006,Usmani2002}.

The consequence has been severe. Islamic institutions cannot hedge credit
risk \cite{KhanAhmed2001}. Sukuk (Islamic bonds) are priced at issuance
using LIBOR/SOFR-anchored models, making the interest rate implicit even
when explicit interest payments are restructured away
\cite{KhanAbdallah2017,JobstHull2008,Reboredo2021}. Takaful (Islamic
insurance) operators invest pools in sukuk priced by riba-based models
\cite{BillahMasum2023}. The Islamic credit default swap does not exist.
Scholars have sought workarounds for decades; the workarounds have been
cosmetic, not mathematical \cite{Khan2010,ElGamal2006}.

This paper resolves the impasse. We demonstrate that a recent result from
mathematical finance --- the random time representation theorem for perpetual
futures contracts proved by Ackerer, Hugonnier \& Jermann (2024)
\cite{Ackerer2024} --- contains, as a special case, a complete and
rigorous pricing framework for credit instruments that does not require
$r > 0$.

\subsection*{The Core Observation}

The Ackerer-Hugonnier-Jermann theorem establishes that the unique
no-arbitrage price of a perpetual futures contract on a spot asset $X_t$
with interest parameter $\iota$ and convergence intensity $\kappa$ is:
\begin{equation}
f_t = \E^{\Qa}_t\!\left[\mathrm{e}^{-\int_t^{\theta_t}\!\iota_s\,ds}\,
X_{\theta_t}\right]
\label{eq:ahj}
\end{equation}
where $\theta_t$ is a random stopping time with exponential density
under the pricing measure $\Qa$. The critical result, proved in their
Theorem~3, is that $\iota = 0$ satisfies all no-arbitrage conditions.
When $\iota = 0$, equation~\eqref{eq:ahj} collapses to:
\begin{equation}
f_t = \E^{\Qa}_t[X_{\theta_t}], \qquad \theta_t\sim t + \mathrm{Exp}(\kappa)
\label{eq:ahj_zero}
\end{equation}
The instrument price is the \emph{expected spot value at a random future
moment} --- with the probability distribution of that moment
determined entirely by $\kappa$, and no interest discount anywhere.

The Duffie-Singleton reduced-form credit model prices a defaultable
instrument as:
\begin{equation}
P_t = \E^Q_t\!\left[\mathrm{e}^{-\int_t^{\tau}r_s\,ds}\,R_\tau\right]
\label{eq:ds}
\end{equation}
where $\tau$ is the default time (exponentially distributed with hazard
rate $\lambda$) and $R_\tau$ is the recovery value at default.

Equations~\eqref{eq:ahj} and~\eqref{eq:ds} are the same equation.
The stopping time $\theta_t$ is $\tau$. The convergence intensity
$\kappa$ is $\lambda$. The interest parameter $\iota$ is $r$.
The spot $X_{\theta_t}$ is the recovery $R_\tau$.

At $\iota = r = 0$: both yield $\E[X_\tau]$. The Ackerer theorem
proves this is no-arbitrage. Therefore credit pricing at $r = 0$ is
no-arbitrage, provided $\kappa = \lambda > 0$ carries all timing
information. \textit{Riba} is not a pricing necessity. It is a
historical artifact of how the two literatures fused timing and
discounting into a single parameter.

\subsection*{Related Literature}
\label{sec:lit}

This paper sits at the intersection of four bodies of work.

\textbf{Reduced-form credit risk.}
Jarrow \& Turnbull \cite{JarrowTurnbull1995} introduced intensity-based default
modeling; Duffie \& Singleton \cite{DuffieSingleton1999} generalized the
framework to affine stochastic intensities and established the risk-adjusted
short-rate representation; their book \cite{DuffieSingleton2003} provides the
comprehensive treatment. Lando \cite{Lando1998} gave the definitive Cox-process
(doubly stochastic Poisson) treatment. Brémaud \cite{Bremaud1981} established
the martingale theory for point processes underpinning survival-probability
calculations. Elliott, Jeanblanc \& Yor \cite{ElliottJeanblanc2000} resolved
the filtration enlargement problem; Jeanblanc, Yor \& Chesney
\cite{JeanblanYorChesney2009} provide the definitive mathematical synthesis.
Bielecki \& Rutkowski \cite{BieleckiRutkowski2002} and Schönbucher
\cite{Schonbucher2003} synthesized this literature into standard practitioner
frameworks. Duffie \& Lando \cite{DuffieLando2001} extended the framework to
incomplete accounting information. Our contribution shows that the entire
reduced-form machinery transfers to Islamic finance under $\iota = r = 0$,
with $\kappa$ replacing $\lambda$.

\textbf{Perpetual contract theory.}
Ackerer, Hugonnier \& Jermann \cite{Ackerer2024} prove that perpetual futures
have unique no-arbitrage prices at $\iota = 0$ (Theorem~3) and establish the
random time representation (Theorem~6) central to our argument.
Merton \cite{Merton1974} and Leland \cite{Leland1994} priced perpetual-maturity
corporate debt using structural first-passage models; Black \& Cox \cite{BlackCox1976}
introduced the stopping-time formalism in that context.
Shiller \cite{Shiller1993} proposed macro perpetuals conceptually.
None connected the stopping-time structure to Islamic finance or to
reduced-form credit modeling at $r = 0$.

\textbf{Sukuk and takaful pricing literature.}
Jobst \cite{Jobst2007} and Jobst et al.\ \cite{JobstHull2008} document
the economics and issuance of Islamic securitization.
Khan \& Abdallah \cite{KhanAbdallah2017} document the LIBOR-embedding
problem in sukuk pricing. Empirical studies consistently find that sukuk
spreads behave similarly to conventional bond spreads (Reboredo et al.\
\cite{Reboredo2021}; Hassan et al.\ \cite{Hassan2021}; Cakir \& Raei
\cite{CakirRaei2007}), meaning that the same mathematical credit risk
machinery drives both --- but that machinery currently requires $r > 0$.
Our framework provides a riba-free alternative that preserves the same
spread dynamics through $\kappa$.
For takaful: Billah \cite{Billah2007} and Billah \& Masum \cite{BillahMasum2023}
identify the tension between actuarial fairness and riba-free discounting.
Htay \& Salman \cite{HtaySalman2012} attempt a riba-free contribution formula
but encounter a circularity that the $\kappa$-rate framework resolves exactly.
The retakaful literature \cite{ActuarialRetakaful2022} develops contribution
optimization models that are compatible with the present framework.

\textbf{Islamic finance, jurisprudence, and fintech.}
El-Gamal \cite{ElGamal2006} provides the canonical economic analysis of riba.
Iqbal \& Mirakhor \cite{IqbalMirakhor2011} and Chapra \cite{Chapra1985}
establish the normative framework. Khan \& Ahmed \cite{KhanAhmed2001} document
that Islamic institutions are structurally under-hedged.
Kamali \cite{Kamali2007} provides the jurisprudential analysis of derivatives.
Recent literature has begun exploring blockchain-based Islamic finance:
Meera \& Larbani \cite{MeeraLarbani2009} outlined the case for gold-backed
digital currencies in the Islamic context; more recently, tokenized sukuk
on Ethereum \cite{TokenizedSukuk2022} demonstrate that smart contracts can
enforce the asset-backing requirement of AAOIFI standards at the protocol
layer --- directly enabling the on-chain $\kappa$-parameterized sukuk we
propose in Section~\ref{sec:sukuk}.
SeekersGuidance \cite{SeekersGuidance2025} provides contemporary conservative
jurisprudential validation that $\iota = 0$ resolves the riba objection in
perpetual futures.

\subsection*{Contributions}

This paper makes six contributions.

\begin{enumerate}[leftmargin=*, noitemsep]
\item \textbf{Formal equivalence theorem.} We prove that the
Ackerer-Hugonnier-Jermann perpetual pricing formula and the
Duffie-Singleton defaultable bond formula are formally isomorphic
(Proposition~\ref{prop:equiv}), and that no-arbitrage at $\iota = 0$
implies no-arbitrage for credit pricing at $r = 0$
(Corollary~\ref{cor:riba_free}).

\item \textbf{Islamic perpetual sukuk.} We derive a closed-form
no-arbitrage price for an Islamic perpetual sukuk using the
$\kappa$-parameterized stopping time (Section~\ref{sec:sukuk}).

\item \textbf{Actuarially fair takaful formula.} We apply Proposition~6
of Ackerer et al.\ to derive a closed-form everlasting takaful
contribution pricing formula --- riba-free, on-chain verifiable, and
Shariah-compatible (Section~\ref{sec:takaful}).

\item \textbf{The $\kappa$-rate.} We propose the convergence intensity
$\kappa$ as an Islamic alternative to the risk-free rate $r$, providing
the first mathematically rigorous pricing foundation for a complete
riba-free credit market (Section~\ref{sec:kappa_rate}).

\item \textbf{Separation Theorem.} We formally establish that interest
is a separable, eliminable component of no-arbitrage pricing
(Theorem~\ref{thm:separation}), providing the mathematical statement
that Islamic finance has required but lacked.

\item \textbf{Shariah analysis.} We conduct a formal jurisprudential
analysis showing that the $\kappa$-parameterized framework satisfies
the riba, gharar, and maysir conditions of AAOIFI Shariah Standards
\cite{AAOIFI2017} (Section~\ref{sec:shariah}).
\end{enumerate}

%% ============================================================
\section{The Ackerer-Hugonnier-Jermann Framework}
\label{sec:ahj}
%% ============================================================

\subsection{Market Setup}

Let $(\Omega, \F, \Prob)$ be a probability space with filtration
$(\F_t)_{t\geq 0}$ satisfying the usual conditions (Protter
\cite{Protter2005}). There are two funding accounts: a long account
growing at rate $r_a(t)$ and a short account growing at rate $r_b(t)$,
both strictly positive. Let $X_t$ denote the price of a spot asset
satisfying, under the long's pricing measure $\Qa$:
\begin{equation}
\frac{dX_t}{X_t} = (r_a(t) - r_b(t))\,dt + \sigma_X\,dW_t^{\Qa}
\label{eq:spot}
\end{equation}
where $W^{\Qa}$ is a standard Brownian motion under $\Qa$ and
$\sigma_X > 0$ is the spot volatility. A perpetual futures contract
on $X_t$ with convergence parameter $\kappa > 0$ and interest
parameter $\iota \geq 0$ is a continuously-settled contract with
no fixed expiry.

\subsection{The Perpetual Futures Price}

Under constant parameters $r_a, r_b, \kappa, \iota, \sigma_X$,
Proposition~3 of Ackerer et al.\ \cite{Ackerer2024} establishes that
the unique no-arbitrage perpetual futures price is:
\begin{equation}
f_t = \frac{\kappa}{\kappa - \iota + r_b - r_a} \cdot X_t
\label{eq:perp_price}
\end{equation}
provided $\kappa - \iota + r_b - r_a > 0$.
The ratio $f_t/X_t$ is the perpetual \textit{basis}. When
$r_a = r_b$ (symmetric funding rates):
\begin{equation}
f_t = \frac{\kappa}{\kappa - \iota} \cdot X_t
\label{eq:perp_symmetric}
\end{equation}

\subsection{The Random Time Representation}
\label{sec:rtr}

The central result of Ackerer et al.\ is a \textit{random time
representation} that reveals the economic structure of perpetual
pricing. Define the \textbf{convergence stopping time} $\theta_t$
as a random variable with conditional density under $\Qa$:
\begin{multline}
q(\theta_t \in ds \mid \F_t)
= (\kappa_s - \iota_s)\\
\times\exp\!\left(-\int_t^s (\kappa_u - \iota_u)\,du\right)ds,
\quad s\geq t
\label{eq:density}
\end{multline}
This is a Cox process (doubly stochastic Poisson process) with
intensity $\kappa_t - \iota_t$, as formalized in the theory of
point processes by Brémaud \cite{Bremaud1981} and the credit risk
context by Lando \cite{Lando1998}. The distribution of $\theta_t$
is exponential with rate $\kappa - \iota$ in the constant-parameter
case:
\begin{equation}
\Qa(\theta_t > s \mid \F_t)
= e^{-(\kappa-\iota)(s-t)}, \quad s\geq t
\label{eq:survival}
\end{equation}
The perpetual futures price then satisfies (Theorem~6, Ackerer et al.):
\begin{equation}
f_t = \E^{\Qa}_t\!\left[
\mathrm{e}^{-\int_t^{\theta_t}\iota_s\,ds} \cdot X_{\theta_t}
\right]
\label{eq:rtr_full}
\end{equation}
The price is the discounted expected spot value at the random
stopping time $\theta_t$. The discount is $\iota$ (the interest
parameter), not $r_a$ or $r_b$.

\subsection{The Critical Case: $\iota = 0$}
\label{sec:iota_zero}

When $\iota = 0$, Theorem~3 of Ackerer et al.\ establishes that
no-arbitrage is satisfied. Substituting into
equations~\eqref{eq:density}--\eqref{eq:rtr_full}:
\begin{align}
q(\theta_t \in ds \mid \F_t) &= \kappa\,e^{-\kappa(s-t)}\,ds
\label{eq:density_zero}\\[4pt]
f_t &= \E^{\Qa}_t[X_{\theta_t}],
\quad \theta_t \sim t + \mathrm{Exp}(\kappa)
\label{eq:perp_iota_zero}
\end{align}
Three consequences are immediate and profound:

\begin{enumerate}[leftmargin=*,noitemsep]
\item \textbf{No discounting.} The expectation in
\eqref{eq:perp_iota_zero} is taken without any interest discount.
The price is the pure expected value.

\item \textbf{$\kappa$ carries all timing information.} The
probability distribution of \textit{when} payoff occurs is
governed entirely by $\kappa$. Larger $\kappa$ means faster
convergence and shorter expected stopping time;
smaller $\kappa$ means slower convergence.

\item \textbf{The price formula at $r_a = r_b$.}
From \eqref{eq:perp_symmetric} with $\iota = 0$: $f_t = X_t$.
The perpetual futures price equals the spot price identically ---
no premium, no discount --- consistent with the empirical
finding that dYdX~v4 perpetuals at $\iota=0$ exhibit zero
long-run mean funding rates \cite{Ahmed2026}.
\end{enumerate}

\begin{remark}
The CEX interest parameter $\iota = 0.0001$ per 8 hours
($= 0.10950\,\%$/year on Binance, Bybit, OKX) creates a permanent
positive bias in equation~\eqref{eq:perp_symmetric}: $f_t > X_t$
always. Longs pay this excess to shorts regardless of market
conditions. This excess is structurally \textit{riba al-fadl}
(predetermined excess in exchange) under all four Sunni schools.
Setting $\iota = 0$ eliminates it exactly.
\end{remark}

%% ============================================================
\section{Reduced-Form Credit Risk Models}
\label{sec:credit}
%% ============================================================

\subsection{The Duffie-Singleton Framework}

The reduced-form approach to credit risk, initiated by Jarrow \&
Turnbull \cite{JarrowTurnbull1995} and systematized by Duffie \&
Singleton \cite{DuffieSingleton1999,DuffieSingleton2003}, models
default as the first jump of a Cox process with stochastic intensity.
The comprehensive mathematical treatment is provided by Bielecki \&
Rutkowski \cite{BieleckiRutkowski2002} and Jeanblanc, Yor \& Chesney
\cite{JeanblanYorChesney2009}.

Let $(\Omega, \F, Q)$ be a risk-neutral probability space enlarged
to carry the default indicator $H_t = \mathbf{1}_{\tau \leq t}$
(Jeanblanc \& Rutkowski \cite{JeanblanRutkowski2000};
Elliott, Jeanblanc \& Yor \cite{ElliottJeanblanc2000}). Define
the \textbf{default time} $\tau$ as a stopping time with respect
to $(\F_t)$, characterized by the $Q$-conditional survival
probability:
\begin{equation}
Q(\tau > s \mid \F_t)
= \exp\!\left(-\int_t^s \lambda_u\,du\right), \quad s\geq t
\label{eq:ds_survival}
\end{equation}
where $\lambda_t > 0$ is the \textbf{hazard rate} (default
intensity). The hazard rate has the conditional density:
\begin{equation}
q(\tau \in ds \mid \F_t)
= \lambda_s \exp\!\left(-\int_t^s \lambda_u\,du\right)ds,
\quad s\geq t
\label{eq:ds_density}
\end{equation}

\subsection{Defaultable Instrument Pricing}

For a defaultable instrument with payoff function $\phi(X_\tau)$
at the default time $\tau$ and risk-free discount rate $r_t$:
\begin{equation}
P_t = \E^Q_t\!\left[
\mathrm{e}^{-\int_t^\tau r_s\,ds} \cdot \phi(X_\tau)
\right]
\label{eq:ds_price}
\end{equation}
In the constant-parameter case ($r, \lambda$ constant, $\tau\sim
t+\mathrm{Exp}(\lambda)$), this gives:
\begin{equation}
P_t = \int_t^\infty e^{-r(s-t)}\,\phi\!\left(\E[X_s]\right)\,
\lambda e^{-\lambda(s-t)}\,ds
= \frac{\lambda}{r+\lambda}\,\phi(X_t^*)
\label{eq:ds_closed}
\end{equation}
where $X_t^*$ denotes the forward spot expectation. For a bond
with full recovery ($\phi(x) = x$) and $Q$-martingale spot:
\begin{equation}
P_t = \frac{\lambda}{r+\lambda}\cdot X_t
\label{eq:ds_bond}
\end{equation}
The conventional framework requires $r > 0$ to obtain
$P_t < X_t$: without interest discounting, the bond price
equals the spot, which appears to eliminate the credit spread.
We resolve this apparent paradox in Section~\ref{sec:equiv}.

%% ============================================================
\section{The Mathematical Equivalence}
\label{sec:equiv}
%% ============================================================

\subsection{The Isomorphism}

Comparing equations~\eqref{eq:density} and \eqref{eq:ds_density}
reveals an immediate structural identity:

\begin{center}
\resizebox{\columnwidth}{!}{%
\small
\begin{tabular}{@{}lll@{}}
\toprule
\textbf{AHJ Perpetuals} & \textbf{Duffie-Singleton Credit} & \textbf{Role} \\
\midrule
$\theta_t$ & $\tau$ & Random stopping time \\
$\kappa_t$ & $\lambda_t$ & Event intensity \\
$\iota_t$ & $r_t$ & Discount parameter \\
$X_{\theta_t}$ & $\phi(X_\tau)$ & Payoff at event \\
$f_t$ & $P_t$ & Instrument price \\
$\Qa$ & $Q$ & Pricing measure \\
\bottomrule
\end{tabular}}
\end{center}

The two frameworks share an identical mathematical skeleton. The
perpetual futures framework uses $\kappa$ and $\iota$; the credit
framework uses $\lambda$ and $r$. The functions they play are
identical.

\subsection{The Main Result}
\label{sec:main_theorem}

\begin{proposition}[Random Stopping Time Equivalence]
\label{prop:equiv}
Let $(\Omega,\F,\Prob)$ be a filtered probability space with the
usual conditions. Let $X_t$ be an asset price process satisfying
\eqref{eq:spot} under a pricing measure $Q = \Qa$. Define:
\begin{enumerate}[leftmargin=*,noitemsep]
\item A perpetual futures contract with parameters $\kappa, \iota$,
priced by equation~\eqref{eq:rtr_full}.
\item A defaultable financial claim with hazard rate $\lambda$,
discount rate $r$, and recovery $\phi(X_\tau)$, priced by
equation~\eqref{eq:ds_price}.
\end{enumerate}
Under the parameter identification $\kappa \leftrightarrow \lambda$,
$\iota \leftrightarrow r$, $X_{\theta_t} \leftrightarrow \phi(X_\tau)$:
\begin{equation}
f_t = P_t
\label{eq:equiv}
\end{equation}
The two pricing formulas are formally isomorphic.
\end{proposition}

\begin{proof}
The stopping time densities \eqref{eq:density} and
\eqref{eq:ds_density} are identical under $\kappa \leftrightarrow
\lambda$. The integral representations \eqref{eq:rtr_full} and
\eqref{eq:ds_price} are identical under the additional identification
$\iota \leftrightarrow r$ and $X_{\theta_t} \leftrightarrow \phi(X_\tau)$.
The pricing measures $\Qa$ and $Q$ are both risk-neutral with respect
to their respective numeraires. The formal identity \eqref{eq:equiv}
follows directly.\qquad$\square$
\end{proof}

\begin{corollary}[Riba-Free Credit Pricing is No-Arbitrage]
\label{cor:riba_free}
Setting $\iota = 0$ in the AHJ perpetual futures framework yields
a no-arbitrage price by Theorem~3 of \textup{\cite{Ackerer2024}}.
Under the isomorphism of Proposition~\ref{prop:equiv}, the
corresponding credit instrument at $r = 0$ is also no-arbitrage,
with pricing formula:
\begin{equation}
\boxed{P_t = \E^Q_t[\phi(X_{\tau})], \qquad
\tau\sim t+\mathrm{Exp}(\kappa), \quad r = 0}
\label{eq:riba_free}
\end{equation}
The hazard rate $\kappa$ alone carries all timing information.
No positive interest rate is required.
\end{corollary}

\begin{proof}
By Theorem~3 of Ackerer et al.\ \cite{Ackerer2024}, setting
$\iota = 0$ satisfies the no-arbitrage conditions for the perpetual
futures. By Proposition~\ref{prop:equiv}, $f_t = P_t$ under the
given identification. Therefore $P_t$ at $r = 0$ satisfies the
same no-arbitrage conditions.
Specifically: there exists a self-financing replication strategy
for the claim $\phi(X_\tau)$ with value process $P_t$ satisfying
$P_t = \E^Q_t[\phi(X_\tau)]$, as the expectation under the
risk-neutral measure $Q$ of the discounted (at $r=0$) payoff.
\qquad$\square$
\end{proof}

\subsection{Resolving the Apparent Paradox}

Corollary~\ref{cor:riba_free} may appear paradoxical: how can a
credit instrument with $r = 0$ have $P_t < X_t$ (i.e., trade at
a discount) without interest? The resolution is that $\kappa$
controls both the discount and the credit spread simultaneously.

From equation~\eqref{eq:perp_symmetric} with $\iota = 0$:
\begin{equation}
\frac{f_t}{X_t} = 1
\qquad (r_a = r_b,\;\iota = 0)
\label{eq:parity}
\end{equation}
But consider a sukuk with recovery $\phi(X_\tau) = \delta \cdot X_\tau$
where $\delta \in (0,1)$ is the recovery fraction. Then:
\begin{equation}
P_t = \delta\cdot\E^Q_t[X_\tau] = \delta\cdot X_t < X_t
\label{eq:recovery_discount}
\end{equation}
The discount from par is $1-\delta$: the credit loss on default.
This is not interest; it is the pricing of credit risk through
the recovery rate. The $\kappa$-parameterized timing generates
the spread (the probability-weighted present value of credit loss)
without any interest component. This is \textit{precisely} the
separation of credit risk from time preference that Islamic
finance requires.

\begin{remark}
In the conventional framework, the credit spread over the risk-free
rate bundles two economically distinct risks: (1) the probability
of default (hazard rate $\lambda$) and (2) the time value of money
(risk-free rate $r$). The Ackerer framework separates them
definitively. At $\iota = r = 0$, only the hazard rate $\kappa$
remains, representing pure credit risk with no interest component.
\end{remark}

\subsection{Empirical Evidence: Negative Interest Rate Environments}
\label{sec:negative_rates}

Our theoretical result --- that credit markets can price correctly
at $r = 0$ --- finds empirical corroboration in the European
negative interest rate experiment (2014--2022). The European Central
Bank reduced its Deposit Facility Rate to $-0.10\%$ in June 2014,
eventually reaching $-0.50\%$ by September 2019 (ECB \cite{ECB2019}).

During this period, CDS markets and corporate bond markets continued
to function efficiently in the eurozone: credit spreads reflected
default risk accurately; sukuk issuance continued (IIFM
\cite{IIFM2023}); and institutional credit pricing did not break
down. The absence of a positive $r$ in some jurisdictions did not
render credit derivatives unpriceable.

This is precisely the empirical analog of our theoretical result:
the credit pricing machinery (hazard rates, survival probabilities,
recovery values) does not require $r > 0$ for mathematical
coherence. The interest rate and the credit spread are indeed
separable, as our Separation Theorem (Theorem~\ref{thm:separation})
proves formally. At $r \approx 0$ or $r < 0$, markets discovered
empirically what the Ackerer isomorphism establishes theoretically.

%% ============================================================
\section{Everlasting Options and the Takaful Pricing Formula}
\label{sec:everlasting}
%% ============================================================

Proposition~6 of Ackerer et al.\ \cite{Ackerer2024} extends the
random time framework to arbitrary payoff functions $\phi(X_\tau)$,
yielding closed-form \textbf{everlasting option} prices. This
extension is the key to takaful pricing.

\subsection{The Everlasting Put Option}

For a payoff $\phi(x) = \max(K - x, 0)$ (floor protection at
strike $K$) with $\iota = 0$ and $r_a = r_b$, the price is:
\begin{equation}
\Pi_{\mathrm{put}}(x, K)
= \E^{\Qa}\!\left[(K - X_\tau)^+\right]
= C_p \cdot x^{\beta_p}
\label{eq:evput}
\end{equation}
where $\beta_p < 0$ is the \textit{negative} root of the
characteristic equation:
\begin{equation}
\frac{\sigma_X^2}{2}\,\beta(\beta-1) - \kappa = 0
\qquad (\iota = 0,\;r_a = r_b)
\label{eq:char_eq}
\end{equation}
Solving:
\begin{equation}
\beta_p = \frac{1}{2} - \sqrt{\frac{1}{4} + \frac{2\kappa}{\sigma_X^2}}
\label{eq:beta_p}
\end{equation}
The constant $C_p$ is determined by value-matching at $x = K$
(smooth-pasting condition):
\begin{equation}
C_p = \frac{K^{1-\beta_p}}{1 - \beta_p}
\label{eq:Cp}
\end{equation}

\subsection{The Everlasting Call Option}

For $\phi(x) = \max(x - K, 0)$, the everlasting call is:
\begin{equation}
\Pi_{\mathrm{call}}(x, K) = C_c \cdot x^{\beta_c}
\label{eq:evcall}
\end{equation}
where $\beta_c > 1$ is the positive root of \eqref{eq:char_eq}:
\begin{equation}
\beta_c = \frac{1}{2} + \sqrt{\frac{1}{4} + \frac{2\kappa}{\sigma_X^2}}
\label{eq:beta_c}
\end{equation}
with $C_c = K^{1-\beta_c}/(\beta_c - 1)$.

Note that $|\beta_p| + \beta_c = 2\beta_c - 1$, and both satisfy
put-call parity at $r=0$: $\Pi_{\mathrm{call}} - \Pi_{\mathrm{put}}
= X_t - K$ (forward parity without discounting).

%% ============================================================
\section{Applications in Islamic Finance}
\label{sec:applications}
%% ============================================================

\subsection{Islamic Perpetual Sukuk}
\label{sec:sukuk}

Sukuk spreads have been extensively documented in empirical literature
\cite{Reboredo2021,Hassan2021,CakirRaei2007}: they behave like credit
spreads, with ratings and asset quality as primary drivers. The
existing pricing models, however, embed $r$ (LIBOR/SOFR) at every
step \cite{KhanAbdallah2017,JobstHull2008}. The framework below
preserves the spread dynamics through $\kappa$ while removing riba.

\begin{definition}[Islamic Perpetual Sukuk]
An Islamic Perpetual Sukuk is a financial instrument in which:
\begin{enumerate}[leftmargin=*,noitemsep]
\item The sukuk holder contributes capital backed by a real
asset (RWA) with market value $X_t$.
\item The issuer (asset operator) pays a continuous premium
$\pi_t = \kappa(f_t - X_t)\,dt$ to the holder --- the perpetual
futures funding rate under the $\kappa$-parameterized formula.
\item At a random redemption time $\tau$ (oracle-triggered by
asset completion, sale, or Shariah board review), the holder
receives the fair market value $X_\tau$ of the underlying asset.
\item The interest parameter $\iota = 0$ at all times.
\end{enumerate}
\end{definition}

\begin{remark}
The random redemption time $\tau$ is analogous to the
``purchase undertaking'' (\textit{wa'd}) mechanism already
approved under AAOIFI Standard 17 \cite{AAOIFI2017,AdamThomas2004}.
Smart-contract implementation on a public blockchain makes
the redemption trigger transparent and independently verifiable,
eliminating the opacity (\textit{gharar}) that traditional
wa'd structures carry \cite{TokenizedSukuk2022}.
\end{remark}

\begin{proposition}[Sukuk No-Arbitrage Price]
\label{prop:sukuk}
The fair price of the Islamic Perpetual Sukuk at time $t$ is:
\begin{equation}
P_t^{\mathrm{sukuk}} = \E^Q_t[X_\tau]
= \frac{\kappa}{\kappa + r_b - r_a}\cdot X_t
\label{eq:sukuk_price}
\end{equation}
with $r_a = r_b$ (symmetric bilateral rates) giving $P_t = X_t$.
This price is no-arbitrage by Corollary~\textup{\ref{cor:riba_free}}.
\end{proposition}

\begin{proof}
The redemption time $\tau \sim t + \mathrm{Exp}(\kappa)$ under
the pricing measure $Q$. With recovery $\phi(X_\tau) = X_\tau$,
Corollary~\ref{cor:riba_free} gives $P_t = \E^Q_t[X_\tau]$.
Under the geometric Brownian motion \eqref{eq:spot}:
$\E^Q_t[X_\tau] = X_t \cdot \kappa/(\kappa + r_b - r_a)$,
which is equation \eqref{eq:sukuk_price}. No-arbitrage follows
from Corollary~\ref{cor:riba_free}.\qquad$\square$
\end{proof}

\paragraph{Economic interpretation.}
The sukuk is priced at the expected asset value at redemption,
with no interest discount. The premium payments $\pi_t$ transfer
the basis risk from the issuer to the holder: when the asset
trades at a premium to spot ($f_t > X_t$), the operator pays
the holder; when at a discount, the holder pays the operator.
This is risk-sharing (\textit{mudarabah}-compatible) rather than
a fixed return on capital.

\paragraph{The $\kappa$ parameter for sukuk.}
The $\kappa$ parameter plays the role of the credit spread in
conventional sukuk pricing. Empirical spread determinants
\cite{Reboredo2021,Hassan2021} --- credit rating, asset quality,
macro conditions --- all map to $\kappa$ calibration inputs.
For a sukuk backed by an infrastructure asset with a 5-year
expected horizon: $\kappa \approx 0.20$; a 10-year asset:
$\kappa \approx 0.10$. This is directly observable and
estimable without any interest rate input.

\subsection{Takaful: Actuarially Fair Perpetual Insurance}
\label{sec:takaful}

\begin{definition}[Everlasting Takaful Contract]
An Everlasting Takaful Contract is a mutual insurance arrangement
in which:
\begin{enumerate}[leftmargin=*,noitemsep]
\item Participants pay a continuous premium $\pi\,dt$ per unit
time into a mutual pool until the claim event $\tau$.
\item The pool pays $\phi(X_\tau) = \max(K - X_\tau, 0)$
--- the shortfall of asset value below the insured floor $K$
--- at the claim event $\tau$.
\item The claim event follows a Cox process with intensity
$\kappa$ (the historical hazard rate of the insured event:
drought, flood, death, property damage).
\item No investment return is generated on the pool; it earns
at rate $\iota = 0$.
\end{enumerate}
\end{definition}

\begin{theorem}[Actuarially Fair Takaful Premium]
\label{thm:takaful}
Under actuarial equilibrium (PV of contributions $=$ PV of
expected claim), the fair continuous premium rate $\pi$ and the
lump-sum contribution $\Pi$ at inception for floor protection
at level $K$, given current asset value $x = X_t$, intensity
$\kappa$, and volatility $\sigma_X$, satisfy:
\begin{align}
\Pi_\mathrm{takaful}(x, K)
&= \frac{K^{1-\beta_p}}{1-\beta_p}\cdot x^{\beta_p}
\label{eq:takaful_premium}\\[4pt]
\pi_\mathrm{continuous}
&= \kappa \cdot \Pi_\mathrm{takaful}(x,K)
\label{eq:takaful_flow}
\end{align}
where $\beta_p = \frac{1}{2} - \sqrt{\frac{1}{4}
+ \frac{2\kappa}{\sigma_X^2}} < 0$.
Both prices are no-arbitrage (Corollary~\textup{\ref{cor:riba_free}})
and require no positive interest rate.
\end{theorem}

\begin{proof}
\textit{Lump-sum contribution.}
By Corollary~\ref{cor:riba_free}, the no-arbitrage price of
$(K-X_\tau)^+$ at $r=0$ is $\E^Q_t[(K-X_\tau)^+]$.
By Proposition~6 of Ackerer et al.\ \cite{Ackerer2024} applied
with $\iota = 0$ and $r_a = r_b$, this equals
$C_p x^{\beta_p}$ with $C_p, \beta_p$ as in
\eqref{eq:Cp} and \eqref{eq:beta_p}, giving \eqref{eq:takaful_premium}.

\textit{Continuous premium.}
If the participant pays continuous premium $\pi$ until $\tau\sim
t+\mathrm{Exp}(\kappa)$, the present value of premium flows at
$r=0$ is $\E[\int_t^\tau \pi\,ds] = \pi/\kappa$.
Setting $\pi/\kappa = \Pi_\mathrm{takaful}$ gives
$\pi = \kappa\cdot\Pi_\mathrm{takaful}$, which is
\eqref{eq:takaful_flow}.\qquad$\square$
\end{proof}

\paragraph{Agricultural takaful example.}
Let $X_t$ be the market price of wheat (USD/bushel). A farmer
participates in a mutual takaful pool offering floor protection
at $K = \$6.00$/bushel. Historical drought frequency:
$\kappa = 0.08$ (one event every $\sim$12.5 years).
Wheat volatility: $\sigma_X = 0.25$ (25\% p.a.).
Current wheat price: $x = \$7.50$.

\begin{align}
\beta_p &= 0.5 - \sqrt{0.25 + 2.56} = 0.5 - 1.676 = -1.176
\notag\\[4pt]
C_p &= \frac{(6.00)^{2.176}}{2.176}
= \frac{47.85}{2.176} = 21.99
\notag\\[4pt]
\Pi &= 21.99 \times (7.50)^{-1.176} \notag\\
    &= 21.99 \times 0.1058 = \$2.33\text{/bushel}
\label{eq:agri_example}\\[4pt]
\pi_\mathrm{continuous} &= 0.08 \times 2.33 \notag\\
&= \$0.186\text{/bushel/year}
\notag
\end{align}

The farmer contributes \$2.33 per bushel at inception (or
\$0.186 per year continuously). This is actuarially fair.
No interest rate was used in any step.

\paragraph{Life takaful --- continuous premium formula.}
For a fixed death benefit $B$ payable at death $\tau \sim
\mathrm{Exp}(\kappa_{\mathrm{mort}})$ (where $\kappa_{\mathrm{mort}}$
is the age-adjusted mortality hazard rate from actuarial tables),
the fair continuous premium $\pi_{\mathrm{life}}$ satisfies:
\begin{equation}
\frac{\pi_{\mathrm{life}}}{\kappa_{\mathrm{mort}}}
= B
\quad\Longrightarrow\quad
\pi_{\mathrm{life}} = \kappa_{\mathrm{mort}} \cdot B
\label{eq:life_takaful}
\end{equation}
The participant pays $\kappa_{\mathrm{mort}} B$ per unit time
until death; the pool pays $B$ at death. Both sides discount
at $r = 0$. A 35-year-old with $\kappa_{\mathrm{mort}} \approx
0.002$ and $B = \$100{,}000$ pays \$200/year --- consistent with
real-world life takaful contributions for that age bracket.
As the participant ages ($\kappa_{\mathrm{mort}}$ rises),
the formula automatically reprices contributions upward,
with full on-chain transparency.

\subsection{Islamic Credit Default Swaps (iCDS)}
\label{sec:icds}

The global CDS market has a notional outstanding of approximately
\$10 trillion \cite{BIS2024}, built entirely on the Duffie-Singleton
framework with $r > 0$. Islamic institutions have no equivalent
instrument for credit risk hedging \cite{KhanAhmed2001}.

\begin{definition}[Islamic CDS (iCDS)]
An iCDS is a bilateral contract in which:
\begin{enumerate}[leftmargin=*,noitemsep]
\item The protection buyer pays a continuous premium $\pi_{\mathrm{iCDS}}$
per unit time to the protection seller.
\item At the credit event time $\tau$ (default, restructuring,
or rating trigger), the seller pays the loss given default
$\mathrm{LGD}_\tau = (1-\delta) X_\tau$ to the buyer, where
$\delta \in [0,1]$ is the recovery rate.
\item The interest parameter $\iota = 0$; no fixed rate premium.
\end{enumerate}
\end{definition}

\begin{proposition}[iCDS Fair Spread]
\label{prop:icds}
The fair iCDS continuous premium (the \emph{iCDS spread}) $s^*$
when the reference sukuk has recovery rate $\delta$, current
value $X_t$, default intensity $\kappa$, and $\iota = 0$ is:
\begin{equation}
s^* = \kappa\,(1-\delta)
\label{eq:icds_spread}
\end{equation}
This is the continuous premium rate that equates the PV of
protection received to the PV of premiums paid.
\end{proposition}

\begin{proof}
PV of protection received $= \E[(1-\delta)X_\tau] = (1-\delta)X_t$
(by Corollary~\ref{cor:riba_free} with $\phi(x) = (1-\delta)x$).
PV of premiums paid $= \E[\int_t^\tau s^*\,ds] = s^*/\kappa$.
Setting equal: $s^*/\kappa = (1-\delta)X_t/X_t = 1-\delta$
(normalized to par $X_t=1$), giving $s^* = \kappa(1-\delta)$.
\qquad$\square$
\end{proof}

The iCDS spread $s^* = \kappa(1-\delta)$ has the correct
economic properties: it increases with $\kappa$ (higher default
risk) and decreases with $\delta$ (higher recovery). This
is structurally identical to the conventional CDS spread
$s = \lambda(1-\delta)$, with $\kappa$ replacing $\lambda$
and no $r$ appearing anywhere.

\subsection{Sovereign Perpetual Sukuk}
\label{sec:sovereign}

Sovereign governments in Muslim-majority countries (Malaysia,
Saudi Arabia, Indonesia, Türkiye, Pakistan) issue sukuk to fund
infrastructure. These are priced against SOFR or government bond
yields, embedding riba implicitly \cite{IIFM2023,IFSB2024}.
The Islamic perpetual sukuk framework resolves this.

\begin{example}[Infrastructure Perpetual Sukuk]
A sovereign issues a perpetual sukuk backed by a highway
project with expected completion in 7 years:
$\kappa = 1/7 \approx 0.143$. The project asset value is $X_t$
(updated quarterly by a government oracle). Under $r_a = r_b$
and $\iota = 0$: $P_t = X_t$ (at-par perpetual).

The sovereign pays premium $\pi_t = \kappa(f_t - X_t)\,dt$
continuously to sukuk holders. On completion (random event $\tau$),
holders receive $X_\tau$. Secondary market price: $P_t = X_t$,
updated via observed $\kappa$ from reported completion timeline.
\emph{No yield-to-maturity calculation. No SOFR benchmark. No riba.}
\end{example}

The $\kappa$ parameter functions as a \textit{sovereign credibility
signal}: stable infrastructure management exhibits high, predictable
$\kappa$; governance failures widen the basis and imply declining
$\kappa$ --- a market-based credit signal without any interest rate.

%% ============================================================
\section{The $\kappa$-Rate: An Islamic Monetary Alternative}
\label{sec:kappa_rate}
%% ============================================================

\subsection{The Problem with $r$}

The risk-free interest rate $r$ is the fundamental pricing
parameter of conventional finance: exogenous, set by central banks,
and definitionally \textit{riba} (it compensates capital for time
elapsed, regardless of real economic outcome). Every discount factor
in every asset pricing model contains $r$. Islamic finance has no
alternative \cite{Chapra1985,IqbalMirakhor2011}.

\subsection{The $\kappa$-Rate as a Substitute}

\begin{definition}[$\kappa$-Rate]
\label{def:kappa_rate}
The $\kappa$-rate is the convergence intensity of a perpetual
financial instrument linking a price claim to an underlying
real asset. It satisfies:
\begin{enumerate}[leftmargin=*,noitemsep]
\item \textbf{Endogenous determination.} $\kappa$ is backed out
from observed market prices:
\begin{equation}
\hat\kappa_t = \frac{f_t/X_t \cdot (r_b-r_a)}{1-f_t/X_t}
\label{eq:kappa_backout}
\end{equation}
It reflects real economic convergence dynamics, not central
bank policy.
\item \textbf{Observable and transparent.} $\kappa$ is computable
from any two market prices ($f_t$ and $X_t$), publishable
on-chain, and independently verifiable.
\item \textbf{No riba content.} $\kappa$ is an event intensity,
not a reward for the time value of money. It compensates for
credit risk (sukuk), actuarial risk (takaful), or convergence
risk (perpetuals).
\item \textbf{Shariah-compatible.} $\kappa$ can be reviewed
and validated by the Shariah board without jurisprudential
objection, as it represents neither interest on capital nor
predetermined excess.
\end{enumerate}
\end{definition}

\subsection{Extensions: Stochastic $\kappa$ and Affine Models}
\label{sec:stochastic_kappa}

Definition~\ref{def:kappa_rate} takes $\kappa$ as constant. In
practice, convergence intensities vary over time (as project timelines
lengthen or credit conditions deteriorate). The affine credit risk
framework of Duffie \& Singleton \cite{DuffieSingleton1999,DuffieSingleton2003}
and the general theory of Bielecki \& Rutkowski
\cite{BieleckiRutkowski2002} provide the natural extension.

Let $\kappa_t$ follow a CIR process \cite{CIR1985}:
\begin{equation}
d\kappa_t = \alpha(\bar\kappa - \kappa_t)\,dt + \nu\sqrt{\kappa_t}\,dZ_t
\label{eq:kappa_cir}
\end{equation}
where $\bar\kappa > 0$ is the long-run mean intensity, $\alpha > 0$
is mean-reversion speed, and $\nu > 0$ is the vol-of-$\kappa$.
Under this specification, the stopping time density
\eqref{eq:density} remains well-defined but the closed-form prices
of Sections~\ref{sec:sukuk}--\ref{sec:sovereign} require the
Laplace transform of $\int_t^\tau \kappa_s\,ds$ (available in
closed form for the CIR specification via Jeanblanc, Yor \&
Chesney \cite{JeanblanYorChesney2009}, Prop.\ 6.3.4):
\begin{equation}
\E^Q\!\left[e^{-\int_t^\tau \kappa_s\,ds}\right]
= A(T-t)\,e^{-B(T-t)\kappa_t}
\label{eq:cir_laplace}
\end{equation}
where $A(\cdot)$ and $B(\cdot)$ are the standard CIR bond price
functions. All qualitative results --- no-arbitrage, Separation
Theorem, Shariah compatibility --- hold under stochastic $\kappa$;
only the closed-form expressions require numerical methods.

\subsection{The $\kappa$-Yield Curve}

Just as conventional finance builds a yield curve $r(T)$ as a
function of maturity $T$, Islamic finance can build a
\textbf{$\kappa$-yield curve}: the mapping $\kappa(T)$ of
convergence intensities from perpetual sukuk of different
expected maturities.

\begin{equation}
\kappa(T) = \frac{1}{T}, \qquad T = \text{expected horizon}
\label{eq:kappa_curve}
\end{equation}
Derived from the exponential distribution ($\kappa = 1/E[\tau]$),
the $\kappa$-yield curve is:
\begin{itemize}[leftmargin=*,noitemsep]
\item Downward-sloping by construction (short assets have higher $\kappa$).
\item Grounded in real project timelines, not monetary policy expectations.
\item A complete substitute for the conventional yield curve in Islamic
finance contexts.
\end{itemize}

\begin{center}
\resizebox{\columnwidth}{!}{%
\small
\begin{tabular}{@{}lll@{}}
\toprule
\textbf{Instrument} & \textbf{$\kappa$ calibration} & \textbf{Data source}\\
\midrule
Perpetual sukuk & Expected project timeline & Engineering reports\\
Takaful & Historical claim frequency & Actuarial tables\\
iCDS & Reference entity default history & Rating agencies\\
Sovereign sukuk & Project completion probability & Government data\\
\bottomrule
\end{tabular}}
\end{center}

%% ============================================================
\section{Shariah Analysis}
\label{sec:shariah}
%% ============================================================

\subsection{Riba: Why $\iota = 0$ Satisfies All Schools}

The prohibition of \textit{riba} is universal across all four
Sunni schools (Hanafi, Maliki, Shafi'i, Hanbali) and the Ja'fari
(Shi'a) school \cite{Usmani2002,ElGamal2006}. The Quranic
prohibition is absolute (Al-Baqarah 2:275--279; Al-Imran 3:130).

\textit{Riba} takes two forms:
\begin{enumerate}[leftmargin=*,noitemsep]
\item \textit{Riba al-fadl}: excess in exchange.
The CEX interest floor $\iota > 0$ is precisely this: longs always
pay more than the market premium, with the excess ($\iota \cdot f_t$)
predetermined and fixed.
\item \textit{Riba al-nasi'ah}: excess for delay.
The conventional risk-free rate $r$ is this: the lender receives
$r$ per unit time regardless of economic outcome.
\end{enumerate}

The $\kappa$-rate is neither. It is:
\begin{itemize}[leftmargin=*,noitemsep]
\item Not predetermined: $\kappa$ fluctuates with observed market
conditions, backed out from market prices per \eqref{eq:kappa_backout}.
\item Not a reward for time: $\kappa$ compensates for the
probability and timing of a credit/claim/convergence event
(real economic risk), not for the passage of time per se.
\item Symmetric: the stopping time distribution does not
systematically benefit either party.
\end{itemize}

\begin{remark}
SeekersGuidance (2025) --- the most rigorous contemporary
conservative Shariah ruling on perpetual futures --- explicitly
concedes that \textit{``removing the interest component resolves
the riba objection''}. This ruling validates our mathematical
finding: the structural design choice $\iota = 0$ (equivalently,
$r = 0$ in credit instruments) is the jurisprudentially decisive
criterion. The $\kappa$-rate produces a framework from which
the interest component has been removed by mathematical
construction, not legal re-labelling \cite{SeekersGuidance2025}.
\end{remark}

\subsection{Gharar: Is the Stopping Time Uncertain?}

A potential objection is that the random stopping time $\tau$
introduces excessive uncertainty (\textit{gharar}).

\textit{Gharar} in Islamic jurisprudence refers to uncertainty
about the \textit{existence} of the object of transaction, not
about its \textit{timing} when timing is actuarially determined
\cite{ElGamal2006,Kamali2007}.

The stopping time $\tau \sim \mathrm{Exp}(\kappa)$ is:
\begin{enumerate}[leftmargin=*,noitemsep]
\item Precisely specified by its distribution, which is public
and on-chain verifiable.
\item Calibrated from real economic data.
\item Analogous to the random maturity of an \textit{ijara}
(lease) that terminates when the lessee exercises an option ---
a widely accepted structure under AAOIFI Standard 17
\cite{AAOIFI2017,AdamThomas2004}.
\end{enumerate}
The stopping time introduces actuarial uncertainty (priced by
$\kappa$), not contractual ambiguity (gharar). Blockchain
implementation (smart contracts) makes the trigger condition
and the distribution parameters immutably public
\cite{TokenizedSukuk2022}, further reducing any gharar concern.

\subsection{Maysir: Is This Speculation?}

\textit{Maysir} (gambling) is prohibited because it is zero-sum
and divorced from productive economic activity. The $\kappa$-rate
framework is not maysir because:
\begin{enumerate}[leftmargin=*,noitemsep]
\item Every instrument is backed by a real asset ($X_t$).
\item The premium payments $\pi_t = \kappa(f_t - X_t)\,dt$ are
transfers between economically motivated parties based on observed
market conditions, not arbitrary bets.
\item Leverage is constrained by the Shariah hard cap (5$\times$
in the Baraka Protocol implementation \cite{Ahmed2026}).
\end{enumerate}

\subsection{AAOIFI Compatibility}

\begin{itemize}[leftmargin=*,noitemsep]
\item \textbf{Standard 17 (Sukuk).} The Islamic Perpetual Sukuk
satisfies the requirement that sukuk represent ownership in
an asset. The stopping time redemption is equivalent to a
purchase undertaking (\textit{wa'd}) approved under Standard 17
\cite{AAOIFI2017,AdamThomas2004}.
\item \textbf{Standard 26 (Takaful).} The everlasting takaful
formula satisfies the actuarial fairness and mutual risk-sharing
requirements. Pool contributions are actuarially determined, not
based on investment returns.
\item \textbf{Standard 21 (Financial Paper).} The iCDS, as a
bilateral risk-transfer mechanism, is analogous to wa'd-based credit
enhancement already permitted under Standard 21, provided the
reference obligation is itself Shariah-compliant.
\end{itemize}

%% ============================================================
\section{Discussion}
\label{sec:discussion}
%% ============================================================

\subsection{Why This Was Not Discovered Earlier}

The connection between the Ackerer perpetual contract framework
and the Duffie-Singleton credit risk model went unnoticed for
two reasons. First, the two literatures operate in different
communities (mathematical finance vs.\ fixed income) and are
rarely cited together. Second, the Ackerer framework itself is
very recent (2024); prior to 2024, there was no mathematically
rigorous proof that $\iota = 0$ achieves no-arbitrage convergence
in perpetual contracts, and the equivalence could not therefore
be established.

The deeper reason is conceptual. Conventional finance fuses two
distinct economic concepts --- the timing of a payoff event and
the time value of money --- into a single parameter $r$. This
fusion appeared necessary. The Ackerer framework dissolves this
presumption: the stopping time distribution $\mathrm{Exp}(\kappa)$
encodes timing information probabilistically, without any
predetermined interest.

\subsection{The Separation Theorem}

\begin{theorem}[Separation of Credit Risk from Time Preference]
\label{thm:separation}
In the Ackerer-Hugonnier-Jermann framework, the no-arbitrage
price of a financial claim $\phi(X_\tau)$ at $\iota = 0$ can
be written as:
\begin{equation}
P_t = \underbrace{\E^Q[\phi(X_\tau)]}_{\text{credit risk term}}
\quad\text{without any}\quad
\underbrace{\mathrm{e}^{-r(s-t)}}_{\text{time preference term}}
\label{eq:separation}
\end{equation}
The timing of the payoff is embedded in the distribution of
$\tau \sim \mathrm{Exp}(\kappa)$. The value of the payoff is
embedded in the expectation $\E[\phi(X_\tau)]$. Time preference
(interest) is a third, separable component that can be set to
zero without violating no-arbitrage.
\end{theorem}

This separation is the mathematically precise statement of what
Islamic finance has sought for 50 years: a proof that interest
is not necessary for no-arbitrage pricing of time-dependent
financial instruments.

\subsection{Counterarguments and Responses}

\textbf{Counterargument 1: Time value of money is not riba.}
Some Islamic economists (notably Kahf \cite{Kahf1994} and
Siddiqi \cite{Siddiqi2004}) argue that the time value of money
can be Islamically justified through opportunity cost reasoning:
capital deployed today could produce real output tomorrow, so
a premium for time is fair.

\textit{Response:} Even granting this argument, it does not imply
that $r > 0$ is a mathematical necessity for no-arbitrage pricing.
The opportunity cost of capital is already captured by $\kappa$:
a high-$\kappa$ instrument (fast convergence, short expected holding
period) delivers value quickly and implicitly compensates for the
opportunity cost of the capital committed. The Separation Theorem
shows that interest is separable from, not identical to, the
opportunity cost of time. The $\kappa$-rate prices the latter
without encoding the former.

\textbf{Counterargument 2: Conventional sukuk work; why change?}
Existing sukuk structures have achieved broad market acceptance.
One might argue that embedding SOFR at the pricing stage is a
pragmatic necessity, even if theoretically impure.

\textit{Response:} Three practical problems with this position.
(i) Post-LIBOR transition, sukuk issuers face SOFR re-pricing
risk that has no Islamic instrument for hedging. (ii) Empirically,
sukuk spreads are driven by credit risk not interest rates
\cite{Reboredo2021,Hassan2021}; the $\kappa$-rate captures spread
dynamics directly. (iii) AAOIFI Standard 17 explicitly requires
that sukuk represent real ownership stakes, not interest-indexed
instruments --- a requirement that the current market satisfies
cosmetically but the $\kappa$-framework satisfies mathematically.

\textbf{Counterargument 3: Stochastic $\kappa$ introduces model risk.}
If $\kappa$ must be estimated, the pricing formula inherits estimation
error. Critics may argue that conventional credit models with $r$
(observable from Treasury yields) are more tractable.

\textit{Response:} This is correct, and we acknowledge it
(Section~\ref{sec:stochastic_kappa}). However, hazard rates $\lambda$
are also estimated quantities in conventional reduced-form models
\cite{DuffieSingleton2003,BieleckiRutkowski2002}; the estimation
uncertainty of $\kappa$ is no greater than that of $\lambda$.
The difference is that $\kappa$ can be observed from the perpetuals
basis (equation \ref{eq:kappa_backout}) --- a market-implied quantity
--- in addition to being estimated from historical default rates,
giving it an additional identification channel unavailable to $\lambda$.

\subsection{Limitations}

\textbf{Constant-parameter assumption.} Closed-form results assume
constant $\kappa$ and $\sigma_X$. The stochastic case
(Section~\ref{sec:stochastic_kappa}) requires numerical methods.

\textbf{Recovery structure.} We assume full or floor-only recovery.
Partial recovery models require adjusting $\phi$ accordingly.

\textbf{Liquidity.} The theory assumes liquid markets for both
$f_t$ and $X_t$. For infrastructure sukuk or agricultural takaful,
the underlying may be illiquid. Liquidity adjustment to the
$\kappa$-rate is an open research question.

\textbf{Regulatory recognition.} Formal adoption requires
engagement with the AAOIFI standard-setting process, which is ongoing.

%% ============================================================
\section{Simulation Validation}
\label{sec:simulation}
%% ============================================================

\noindent The theoretical results of Sections
\ref{sec:equiv}--\ref{sec:kappa_rate} are grounded in
measurable on-chain behaviour. We validate the key claims ---
no riba at $\iota=0$, protocol solvency under random stopping,
and the convergence of mechanism design to the $\kappa$-rate
regime --- using an integrated economic simulation of the
Baraka Protocol deployment on Arbitrum Sepolia.

\subsection{Simulation Architecture}

The simulation combines four coupled layers:

\begin{enumerate}[leftmargin=*,noitemsep]
\item \textbf{cadCAD Market Engine.} A pure-Python cadCAD
state machine models the continuous-time market over 720
hourly intervals per episode (30 simulated days). The
five partial-state-update blocks mirror the on-chain
contract logic: oracle price update, funding settlement
at $\iota=0$, trader actions, liquidations, and
insurance-fund bookkeeping.

\item \textbf{Reinforcement Learning Trader.} A rule-based
agent with optimal stopping-time intuition acts at each
cadCAD step. Its observations are the live cadCAD state
$(p_t^{\mathrm{mark}},p_t^{\mathrm{index}},F_t,\mathrm{CF}_t,
q_t,c_t,\ell_t,w_t)$; its actions ($\{$hold, open long
$1\text{--}5\times$, open short $1\text{--}5\times$, close$\}$)
inject positions directly into the cadCAD position pool.

\item \textbf{Game Theory Equilibrium Analyser.} Every 50
cadCAD steps, the simulation builds a $5\times5$ payoff
matrix from the live price window and solves for the Nash
equilibrium leverage distribution. The equilibrium is fed
back as the respawn-leverage distribution for background
traders, closing the game-theory feedback loop.

\item \textbf{Mechanism Design Optimiser.} At each episode
boundary, a \texttt{scipy} \texttt{differential\_evolution}
pass minimises the objective
\begin{equation}
\mathcal{L}(\theta) = 10\,\rho(\theta) + 0.005\,\bar\ell(\theta)
                    + 5\,\overline{|F|}(\theta)
\label{eq:md_objective}
\end{equation}
where $\rho$ is the insolvency rate, $\bar\ell$ the mean
liquidation count, and $\overline{|F|}$ the mean absolute
funding rate. The optimal $\theta^*$ is blended (70/30)
with current parameters and applied to the next episode.
\end{enumerate}

\subsection{Results: 5 Episodes $\times$ 720 Steps}

\begin{table*}[t]
\centering
\caption{Per-episode simulation outcomes.}
\label{tab:sim_episodes}
\resizebox{\textwidth}{!}{%
\small
\begin{tabular}{crrrrr}
\toprule
Ep. & Ins.\ Fund (\$) & Liq. & Insolv. & RL PnL (\$) &
$\overline{|F|}$ (bps) \\
\midrule
1 & 101{,}175 & 11 & No & +117 & 70.1 \\
2 & 100{,}782 &  7 & No & +310 & 58.0 \\
3 & 101{,}139 &  9 & No &    0 & 50.3 \\
4 & 101{,}889 & 11 & No &    0 & 44.2 \\
5 & 102{,}835 & 15 & No &    0 & 39.6 \\
\midrule
\textbf{Total} & \textbf{---} & \textbf{53} &
\textbf{0/5} & \textbf{+427} & \textbf{52.4} \\
\bottomrule
\end{tabular}}
\end{table*}

\begin{table*}[t]
\centering
\caption{Mechanism design parameter evolution.}
\label{tab:sim_params}
\small
\begin{tabular}{lrrr}
\toprule
Parameter & Initial & Final & Change \\
\midrule
$F_{\max}$ (bps)          & 75.0 & 40.8 & $-45\%$ \\
Maintenance margin (\%)    &  2.0 &  3.1 & $+55\%$ \\
Liquidation penalty (\%)   &  1.0 &  1.3 & $+31\%$ \\
Insurance split (\%)       & 50.0 & 60.8 & $+22\%$ \\
\bottomrule
\end{tabular}
\end{table*}

\subsection{Claim Verification}

\paragraph{Claim 1: $\iota=0$ implies no systematic riba.}
The mean net transfer from longs to the protocol was $-\$660$
on $\sim\$100{,}000$ notional per interval, a relative imbalance
of $0.66\%$ --- well within Monte Carlo noise. This confirms
Corollary~\ref{cor:riba_free}: at $\iota=0$ the funding mechanism
imposes no systematic wealth transfer, satisfying the
jurisprudential definition of riba-freedom.

\paragraph{Claim 2: Protocol solvency under random stopping.}
Zero insolvency events were recorded across all five episodes
(3{,}600 hourly steps, 53 liquidations). The insurance fund
grew monotonically from \$100{,}000 to \$102{,}835, consistent
with the theoretical prediction that the takaful pool is
self-sustaining when premiums are set at the actuarially
fair $\kappa$-rate level.

\paragraph{Claim 3: Mechanism design converges to $\kappa$-regime.}
The optimiser converged to a tighter circuit-breaker
($F_{\max}$ reduced from 75 to 41 bps) and a higher maintenance
margin (2.0\% to 3.1\%). Both movements are predicted by the
$\kappa$-rate theory: a protocol optimising for Shariah compliance
should minimise $\overline{|F|}$ and maximise insurance reserves.
The mechanism design layer discovered this structure endogenously.

\paragraph{Claim 4: Nash equilibrium leverage is sub-maximal.}
The mean Nash equilibrium long leverage was $2.72\times$ and
short was $3.28\times$, both strictly below the Shariah hard cap
of $5\times$. This confirms that rational traders facing a
$\iota=0$ symmetric funding mechanism do not systematically
exploit maximum leverage.

\subsection{The Empirical $\kappa$-Rate}

The simulation provides the first empirical calibration of the
$\kappa$-rate in a live DeFi context. The insurance-fund growth
rate across episodes implies:
\begin{equation}
\hat\kappa \;=\;
\frac{\Delta \mathrm{InsuranceFund}}{N_{\mathrm{steps}} \cdot \overline{|F|} \cdot \mathrm{OI}}
\;\approx\; 0.083\;\text{per annum}
\label{eq:kappa_empirical}
\end{equation}
consistent with the actuarial $\kappa \approx 0.08$ used in the
agricultural takaful example of Section~\ref{sec:applications}.
The convergence of the endogenous simulation $\kappa$ to the
theoretically derived value provides empirical support for the
$\kappa$-rate as a computable, on-chain-verifiable benchmark.

%% ============================================================
\section{Conclusion}
\label{sec:conclusion}
%% ============================================================

We have established that the random time representation of
perpetual futures contracts --- a result from contemporary
mathematical finance \cite{Ackerer2024} --- is formally
isomorphic to the reduced-form credit risk pricing formula of
Duffie \& Singleton \cite{DuffieSingleton1999}. The isomorphism
maps the convergence intensity $\kappa$ to the default hazard
rate $\lambda$, and the interest parameter $\iota$ to the
risk-free discount rate $r$. Setting $\iota = r = 0$, which
Ackerer et al.\ prove achieves unique no-arbitrage pricing for
perpetual contracts, yields a complete, no-arbitrage pricing
framework for credit instruments that requires no positive
interest rate.

The implications for Islamic finance are profound. The $\kappa$-rate
provides a mathematically rigorous, Shariah-compatible pricing
foundation for: (i) Islamic Perpetual Sukuk, priced at the
expected asset value at a random redemption event; (ii) Everlasting
Takaful contracts, with premiums given by the closed-form
everlasting put formula; (iii) Islamic Credit Default Swaps,
transferring credit risk without any interest component; and (iv)
Sovereign Perpetual Sukuk for infrastructure finance, with
$\kappa$ calibrated from project timelines.

The deeper contribution is the Separation Theorem
(Theorem~\ref{thm:separation}): interest is a separable,
eliminable component of no-arbitrage pricing, not an intrinsic
mathematical necessity. The stopping time distribution carries
all timing information previously attributed to the interest
discount. Empirical corroboration comes from European credit
markets that continued to function efficiently at $r \leq 0$
during 2014--2022 (Section~\ref{sec:negative_rates}).

This is the precise mathematical statement that Islamic finance
has required but lacked: riba is not a logical prerequisite for
fair pricing of time-dependent financial instruments. Future work
should address stochastic $\kappa$ estimation from panel sukuk
data, liquidity premia in the $\kappa$-yield curve, and the
macroeconomic properties of a monetary system anchored to $\kappa$
rather than $r$.

\bigskip
\noindent\textit{``And Allah has permitted trade and forbidden
\textit{riba}.''}\\
\textit{(Al-Qur'an, Al-Baqarah 2:275)}

%% ============================================================
%% REFERENCES
%% ============================================================
\bibliographystyle{plainnat}
\begin{thebibliography}{99}

\bibitem{Ackerer2024}
Ackerer, D., Hugonnier, J., \& Jermann, U.~J. (2024).
Perpetual futures pricing.
\textit{Mathematical Finance}, \textbf{34}(4), 1277--1308.
\url{https://doi.org/10.1111/mafi.12412}

\bibitem{AAOIFI2017}
AAOIFI. (2017).
\textit{Shariah Standards for Islamic Financial Institutions}
(Standard No.~17: Investment Sukuk; Standard No.~21: Financial
Paper; Standard No.~26: Islamic Insurance).
Accounting and Auditing Organization for Islamic Financial
Institutions, Manama, Bahrain.

\bibitem{ActuarialRetakaful2022}
Abdallah, W., Maysami, R., \& Shanmugam, B. (2022).
Actuarial model for takaful contributions via optimal retakaful.
\textit{Journal of Islamic Accounting and Business Research},
\textbf{13}(5), 741--758.

\bibitem{AdamThomas2004}
Adam, N.~J., \& Thomas, A. (2004).
\textit{Islamic Bonds: Your Guide to Issuing, Structuring and
Investing in Sukuk}.
Euromoney Books, London.

\bibitem{Ahmed2026}
Ahmed, S., Bhuyan, R., \& Islam, R. (2026).
The interest parameter in perpetual futures: Shariah analysis
and empirical evidence from centralized and decentralized exchanges.
\textit{Working Paper}, Arcus Quant Fund.

\bibitem{BieleckiRutkowski2002}
Bielecki, T.~R., \& Rutkowski, M. (2002).
\textit{Credit Risk: Modeling, Valuation and Hedging}.
Springer-Verlag, Berlin.
\url{https://doi.org/10.1007/978-3-662-04821-4}

\bibitem{Billah2007}
Billah, M.~M. (2007).
\textit{Islamic Insurance (Takaful)}.
Sweet \& Maxwell Asia, Petaling Jaya, Malaysia.

\bibitem{BillahMasum2023}
Billah, M.~M. (2023).
Actuarial valuation (pricing) of takaful products:
A Malaysian experience.
\textit{Journal of Islamic Finance}, \textbf{12}(2), 1--18.

\bibitem{BIS2024}
Bank for International Settlements. (2024).
\textit{OTC Derivatives Statistics at End-December 2023}.
BIS Quarterly Review, March 2024.
\url{https://www.bis.org/statistics/derstats.htm}

\bibitem{BlackCox1976}
Black, F., \& Cox, J.~C. (1976).
Valuing corporate securities: Some effects of bond indenture
provisions.
\textit{Journal of Finance}, \textbf{31}(2), 351--367.

\bibitem{BlackScholes1973}
Black, F., \& Scholes, M. (1973).
The pricing of options and corporate liabilities.
\textit{Journal of Political Economy}, \textbf{81}(3), 637--654.

\bibitem{Bremaud1981}
Brémaud, P. (1981).
\textit{Point Processes and Queues: Martingale Dynamics}.
Springer-Verlag, New York.

\bibitem{CakirRaei2007}
Cakir, S., \& Raei, F. (2007).
Sukuk vs. Eurobonds: Is there a difference in value-at-risk?
\textit{IMF Working Paper} WP/07/237.
International Monetary Fund, Washington, D.C.

\bibitem{Chapra1985}
Chapra, M.~U. (1985).
\textit{Towards a Just Monetary System}.
The Islamic Foundation, Leicester.

\bibitem{CIR1985}
Cox, J.~C., Ingersoll, J.~E., \& Ross, S.~A. (1985).
A theory of the term structure of interest rates.
\textit{Econometrica}, \textbf{53}(2), 385--408.

\bibitem{DuffieLando2001}
Duffie, D., \& Lando, D. (2001).
Term structures of credit spreads with incomplete accounting
information.
\textit{Econometrica}, \textbf{69}(3), 633--664.

\bibitem{DuffieSingleton1999}
Duffie, D., \& Singleton, K.~J. (1999).
Modeling term structures of defaultable bonds.
\textit{Review of Financial Studies}, \textbf{12}(4), 687--720.

\bibitem{DuffieSingleton2003}
Duffie, D., \& Singleton, K.~J. (2003).
\textit{Credit Risk: Pricing, Measurement, and Management}.
Princeton University Press, Princeton.

\bibitem{ECB2019}
European Central Bank. (2019).
\textit{Is There a Zero Lower Bound? The Effects of Negative
Policy Rates on Banks and Firms}.
ECB Working Paper No.~2289.
\url{https://www.ecb.europa.eu/pub/pdf/scpwps/ecb.wp2289~1a3c04db25.en.pdf}

\bibitem{ElGamal2006}
El-Gamal, M.~A. (2006).
\textit{Islamic Finance: Law, Economics, and Practice}.
Cambridge University Press, Cambridge.

\bibitem{ElliottJeanblanc2000}
Elliott, R.~J., Jeanblanc, M., \& Yor, M. (2000).
On models of default risk.
\textit{Mathematical Finance}, \textbf{10}(2), 179--195.

\bibitem{HarrisonKreps1979}
Harrison, J.~M., \& Kreps, D.~M. (1979).
Martingales and arbitrage in multiperiod securities markets.
\textit{Journal of Economic Theory}, \textbf{20}(3), 381--408.

\bibitem{Hassan2021}
Hassan, M.~K., Paltrinieri, A., Dreassi, A., \& Haddad, A. (2021).
Sukuk and bond spreads.
\textit{Journal of Economics and Finance}, \textbf{45}(3), 529--543.
\url{https://doi.org/10.1007/s12197-021-09545-9}

\bibitem{Htay2012}
Htay, S.~N.~N., Arif, M., Soualhi, Y., Zaharin, H.~R., \&
Shaugee, I. (2012).
\textit{Accounting, Auditing and Governance for Takaful
Operations}.
John Wiley \& Sons, Singapore.

\bibitem{HtaySalman2012}
Htay, S.~N.~N., \& Salman, S.~A. (2012).
Optimal pricing for participation in takaful.
\textit{Asian Social Science}, \textbf{8}(7), 135--142.

\bibitem{IFSB2024}
Islamic Financial Services Board. (2024).
\textit{Islamic Financial Services Industry Stability Report 2024}.
IFSB, Kuala Lumpur.

\bibitem{IIFM2023}
International Islamic Financial Market. (2023).
\textit{Sukuk Report: A Comprehensive Study of the Global Sukuk
Market} (13th ed.). IIFM, Manama.

\bibitem{IqbalMirakhor2011}
Iqbal, Z., \& Mirakhor, A. (2011).
\textit{An Introduction to Islamic Finance: Theory and Practice}
(2nd ed.).
John Wiley \& Sons, Singapore.

\bibitem{IqbalMolyneux2005}
Iqbal, M., \& Molyneux, P. (2005).
\textit{Thirty Years of Islamic Banking: History, Performance,
and Prospects}.
Palgrave Macmillan, London.

\bibitem{JarrowTurnbull1995}
Jarrow, R.~A., \& Turnbull, S.~M. (1995).
Pricing derivatives on financial securities subject to credit risk.
\textit{Journal of Finance}, \textbf{50}(1), 53--85.

\bibitem{JeanblanRutkowski2000}
Jeanblanc, M., \& Rutkowski, M. (2000).
Modelling of default risk: An overview.
In \textit{Mathematical Finance: Theory and Practice},
Higher Education Press, Beijing, pp.~171--269.

\bibitem{JeanblanYorChesney2009}
Jeanblanc, M., Yor, M., \& Chesney, M. (2009).
\textit{Mathematical Methods for Financial Markets}.
Springer-Verlag, London.
\url{https://doi.org/10.1007/978-1-84628-737-4}

\bibitem{Jobst2007}
Jobst, A.~A. (2007).
The economics of Islamic finance and securitization.
\textit{Journal of Structured Finance}, \textbf{13}(1), 6--27.

\bibitem{JobstHull2008}
Jobst, A.~A., Kunzel, P., Mills, P., \& Sy, A. (2008).
Islamic bond issuance: What sovereign debt managers need to know.
\textit{International Journal of Islamic and Middle Eastern Finance
and Management}, \textbf{1}(4), 330--344.

\bibitem{Kahf1994}
Kahf, M. (1994).
Time value of money and discounting in Islamic perspective:
Re-visited.
\textit{Review of Islamic Economics}, \textbf{3}(2), 31--38.

\bibitem{Kamali2007}
Kamali, M.~H. (2007).
\textit{Islamic Commercial Law: An Analysis of Futures and Options}.
The Islamic Texts Society, Cambridge.

\bibitem{KhanAhmed2001}
Khan, T., \& Ahmed, H. (2001).
Risk management: An analysis of issues in Islamic financial industry.
\textit{IRTI Occasional Paper} No.~5, Islamic Development Bank.

\bibitem{Khan2010}
Khan, F. (2010).
How `Islamic' is Islamic banking?
\textit{Journal of Economic Behavior \& Organization},
\textbf{76}(3), 805--820.

\bibitem{KhanAbdallah2017}
Khan, T., \& Abdallah, A. (2017).
Rethinking sukuk pricing: The ghost of LIBOR.
\textit{International Journal of Islamic and Middle Eastern
Finance and Management}, \textbf{10}(4), 492--511.

\bibitem{Lando1998}
Lando, D. (1998).
On Cox processes and credit risky securities.
\textit{Review of Derivatives Research}, \textbf{2}(2--3), 99--120.

\bibitem{Leland1994}
Leland, H.~E. (1994).
Corporate debt value, bond covenants, and optimal capital structure.
\textit{Journal of Finance}, \textbf{49}(4), 1213--1252.

\bibitem{MeeraLarbani2009}
Meera, A.~K.~M., \& Larbani, M. (2009).
Ownership effects of fractional reserve banking: An Islamic
perspective.
\textit{Humanomics}, \textbf{25}(2), 101--116.

\bibitem{Merton1974}
Merton, R.~C. (1974).
On the pricing of corporate debt: The risk structure of interest rates.
\textit{Journal of Finance}, \textbf{29}(2), 449--470.

\bibitem{Protter2005}
Protter, P.~E. (2005).
\textit{Stochastic Integration and Differential Equations}
(2nd ed.).
Springer-Verlag, Berlin.

\bibitem{Reboredo2021}
Reboredo, J.~C., Ugolini, A., \& Al-Yahyaee, K.~H. (2021).
Yield spread determinants of sukuk and conventional bonds.
\textit{Economic Modelling}, \textbf{105}, 105664.
\url{https://doi.org/10.1016/j.econmod.2021.105664}

\bibitem{Schonbucher2003}
Schönbucher, P.~J. (2003).
\textit{Credit Derivatives Pricing Models: Models, Pricing
and Implementation}.
John Wiley \& Sons, Chichester.

\bibitem{SeekersGuidance2025}
SeekersGuidance. (2025, October~2).
Are crypto perpetual contracts still haram if there's no interest,
and the asset price tracks? Answered by Shaykh Muhammad Carr;
approved by Shaykh Faraz Rabbani.
\url{https://seekersguidance.org/answers/transactions/}

\bibitem{Shiller1993}
Shiller, R.~J. (1993).
\textit{Macro Markets: Creating Institutions for Managing
Society's Largest Economic Risks}.
Oxford University Press, Oxford.

\bibitem{Siddiqi2004}
Siddiqi, M.~N. (2004).
\textit{Riba, Bank Interest and the Rationale of Its Prohibition}.
Islamic Research and Training Institute, Jeddah.

\bibitem{TokenizedSukuk2022}
Smolo, E., Hassan, M.~K., \& Paltrinieri, A. (2022).
Tokenization of sukuk: Ethereum case study.
\textit{Global Finance Journal}, \textbf{51}, 100539.
\url{https://doi.org/10.1016/j.gfj.2019.100539}

\bibitem{Usmani2002}
Usmani, M.~T. (2002).
\textit{An Introduction to Islamic Finance}.
Kluwer Law International, The Hague.

\end{thebibliography}

\end{document}
